\section{Security Model and Theorems}
\label{sec:security}

This section formalises the adversary model, lists the cryptographic and
system assumptions on which TARDUS rests, and states the six main theorems
together with reduction sketches. Full proof mechanisation is deferred to
Phase 4 (EasyCrypt); the proof sketches given here are sufficient to guide
that work and to serve as the audit baseline.

\subsection{Threat Model}
\label{sec:threat-model}

\paragraph{Adversary classes.} TARDUS considers four overlapping adversary
classes; each theorem in this section specifies which subset it tolerates.

\begin{description}[leftmargin=2em,style=nextline]
  \item[$\Adv_{\mathsf{net}}$ (passive observer)]
    Sees every on-chain transaction, all relay traffic between the user and
    the mint committee, and the contents of the public registries
    (\S5). Cannot inject, drop, or alter messages. This is
    the strictly weakest model.
  \item[$\Adv_{\mathsf{user}}$ (malicious user)]
    A polynomial-time party that controls one or more user identities, can
    interact with the mint committee through the canonical protocols
    (\S3 and \S4), and can submit arbitrary
    transactions to the on-chain program. Adaptive: may corrupt mint
    operators after seeing earlier transcripts.
  \item[$\Adv_{\mathsf{mint},t-1}$ (corrupted mint minority)]
    Statically corrupts up to $t-1$ of the $n$ mint operators (where $t$ is
    the signing threshold). Sees all corrupted operators' shares, internal
    state, and HSM-mediated traffic. Honest operators communicate over
    authenticated channels with replay protection. This is the maximum
    corruption tolerated by all positive results below; at $t$ corruptions
    coin forgery becomes possible by reconstruction.
  \item[$\Adv_{\mathsf{chain}}$ (network-level adversary)]
    Controls the Solana network up to but not including consensus
    (cf.\ Solana's $1/3$ honest-stake assumption). Can reorder
    transactions within a slot, drop unsubmitted transactions, and observe
    all account contents. Cannot subvert the on-chain program's atomic
    state transitions once a transaction is included in a finalised block.
\end{description}

\paragraph{Composability.} Privacy results in this section hold against the
union $\Adv_{\mathsf{net}} \cup \Adv_{\mathsf{mint},t-1} \cup
\Adv_{\mathsf{chain}}$ jointly; double-spend resistance holds additionally
against $\Adv_{\mathsf{user}}$.

\paragraph{What is \emph{not} modelled.} We exclude side-channel attacks on
the user's wallet device (covered by wallet-level mitigations in
\S6), targeted denial-of-service against individual mint
operators (covered by the $t$-out-of-$n$ availability margin in
\S8), and physical compromise of HSMs (covered by the
mandatory FIPS 140-2 Level 3 deployment requirement in \S3.1).

\subsection{Security Assumptions}
\label{sec:security-assumptions}

\paragraph{A1 (ECDLP on $\mathbb{G}$).}
\label{assn:ecdlp}
For the ed25519 prime-order subgroup $\mathbb{G}$ of order
$\ell \approx 2^{252}$, given a uniformly random $X = x \cdot G$, no
polynomial-time adversary recovers $x$ with non-negligible probability:
\[
  \Pr_{x \gets \mathbb{Z}_\ell}\bigl[\Adv(X) = x\bigr] \leq \negl(\lambda).
\]

\paragraph{A2 (Random-Oracle Model).}
\label{assn:rom}
We model SHA-512-reduced-mod-$\ell$ (used for Schnorr challenges,
\S2.3) and SHA-256 (used for nullifiers,
\S4.2) as random oracles in the security proofs of
T1, T2, T4. The hash-to-scalar output bias is bounded by $2^{-260}$
(\S2.3), which is sub-negligible for any
$\poly(\lambda)$ adversary.

\paragraph{A3 (Pedersen binding under DDH on $(G, H)$).}
\label{assn:pedersen}
For the joint commitment $C_{i,0} = s_i \cdot G + r_i \cdot H$
(\S3.5) with independently sampled generator $H$
(canonical hash-to-curve of the domain string \texttt{"TARDUS-vss-h-v1"}),
no polynomial-time adversary opens $C$ to two distinct pairs
$(s_i, r_i), (s_i', r_i')$ with $s_i \neq s_i'$.

\paragraph{A4 (Hash collision resistance).}
\label{assn:cr}
SHA-256 and SHA-512 are collision-resistant against $\poly(\lambda)$
adversaries with output-bit-length security.

\paragraph{A5 (Honest threshold).}
\label{assn:threshold}
In any DKG, sign, or refresh ceremony, at least $t$ of the $n$ mint
operators follow the protocol honestly. Corruption is static
(decided before the ceremony begins); the proactive reshare protocol
(\S3.7) re-randomises shares between epochs to mitigate
slow-creep corruption.

\paragraph{A6 (Solana consensus safety).}
\label{assn:solana}
Transactions included in a Solana block finalised at depth $\geq 32$ slots
are durable; the on-chain program's atomic state transitions hold under
this finality assumption.

\subsection{Theorem T1 --- Coin Unforgeability}
\label{sec:thm-unforgeability}

\begin{theorem}[Coin Unforgeability]
\label{thm:unforgeability}
Under assumptions A1 (ECDLP), A2 (ROM), and A5 (honest threshold), the
TARDUS threshold blind Schnorr signature (\S4) is
existentially unforgeable under chosen-message attack
($\mathsf{EUF\text{-}CMA}$) against any adversary $\Adv \in
\Adv_{\mathsf{user}} \cup \Adv_{\mathsf{mint},t-1}$. Formally, for any
$\poly(\lambda)$-time $\Adv$ making $q_s$ signing queries and $q_h$
random-oracle queries,
\[
  \Pr\bigl[\AdvG{\mathsf{unforge}}^{\Adv}(\lambda) = 1\bigr] \leq
  \frac{q_h \cdot \mathsf{Adv}^{\mathsf{ECDLP}}_{\mathbb{G}}(\lambda)}{1 - 1/\ell}
  + \negl(\lambda).
\]
\end{theorem}

\paragraph{Reduction sketch.} The proof composes two classical results:

\begin{enumerate}[leftmargin=2em,itemsep=0pt]
  \item \emph{Single-key Schnorr is EUF-CMA under ECDLP + ROM} by
        Pointcheval-Stern (\cite{pointcheval-stern1996}) via the
        forking-lemma rewinding argument.
  \item \emph{Threshold lift preserves unforgeability} by Stinson-Strobl
        (\cite{stinson-strobl2001}) when the underlying blind signature is
        unforgeable and the threshold is at most $t-1$. The reduction
        simulates the corrupted minority's view from the random
        oracle programming, exactly as in single-key.
\end{enumerate}

The validator-side commitment phase
(\S4.4) prevents the bias attack of
\cite{nick-ruffing-seurin2020}: validators publish $R_i = k_i \cdot G$
\emph{before} seeing the user's message, so the adversary cannot adaptively
choose challenges to favour a target signature.

\paragraph{Corruption boundary.} At exactly $t$ corruptions the adversary
reconstructs the joint secret $x = \sum_{i \in S} \lambda_{S,i} \cdot s_i$
via Lagrange interpolation (\S3.6) and forgery is trivial.
The proactive reshare protocol (\S3.7) bounds the corruption
window to one epoch.

\subsection{Theorem T2 --- Issuance Blindness}
\label{sec:thm-blindness}

\begin{theorem}[Issuance Blindness]
\label{thm:issuance-blindness}
Under assumption A1 (ECDLP) and the Pointcheval-Stern blinding analysis
(\cite{pointcheval-stern1996}), the user's view in a TARDUS issuance
session is statistically independent of the resulting coin's pair
$(Cp, \sigma)$. Formally, for any unbounded mint-side adversary
$\Adv \in \Adv_{\mathsf{mint},t-1}$ and any two messages
$Cp_0, Cp_1 \in \mathbb{G}$, the issuance transcripts produced by signing
$Cp_0$ and $Cp_1$ are perfectly indistinguishable:
\[
  \mathsf{SD}\bigl(\mathsf{Trans}(Cp_0),\ \mathsf{Trans}(Cp_1)\bigr) = 0.
\]
\end{theorem}

\paragraph{Reduction sketch.} The user samples uniformly random blinding
factors $\alpha, \beta \gets \mathbb{Z}_\ell^\ast$ in Round 2 of the
threshold blind sign (\S4.5). The blinded challenge
$c' = c + \beta$ and blinded commitment $R' = R + \alpha \cdot G + \beta
\cdot Cp$ are uniform over $\mathbb{Z}_\ell$ and $\mathbb{G}$
respectively, independent of $Cp$. The mint sees only $(R', c', s')$
where $s'$ is the response to $c'$; the unblinding step
($s = s' - \alpha$, $c$ recovered) preserves uniformity. The
information-theoretic argument gives statistical distance $0$, stronger
than the computational indistinguishability typically achieved by
zero-knowledge schemes.

\paragraph{Note.} The blindness result is \emph{unconditional} (no
computational assumption on the mint side), in contrast to T1 which
reduces to ECDLP. This is a strength inherited from Chaum's 1982
construction.

\subsection{Theorem T3 --- Double-Spend Prevention}
\label{sec:thm-double-spend}

\begin{theorem}[Double-Spend Prevention]
\label{thm:double-spend}
Under assumptions A4 (SHA-256 collision resistance) and A6 (Solana
finality), any attempt by $\Adv \in \Adv_{\mathsf{user}}$ to submit a
spend for the same coin $(Cp, \sigma)$ twice is rejected with probability
$1$ at the protocol level. Formally,
\[
  \Pr\bigl[\AdvG{\mathsf{double\text{-}spend}}^{\Adv}(\lambda) = 1\bigr] = 0
\]
after the first spend's transaction has finalised on chain.
\end{theorem}

\paragraph{Proof.} The nullifier
$\mathsf{null}(Cp) = \mathsf{SHA\text{-}256}(\text{``TARDUS-nullifier-v1''} \,\|\, Cp)$
is deterministic in $Cp$ (\S4.2). The on-chain
\texttt{Refresh} and \texttt{Withdraw} handlers (\S5.3.3)
perform $\mathsf{NullifierSet.insert}(\mathsf{null}(Cp))$, which returns
\texttt{false} if the entry is already present and aborts the
transaction with \texttt{ERR\_DOUBLE\_SPEND}. This insertion is wrapped
inside the atomic boundary of a single Solana transaction; under A6 the
state mutation is durable upon finalisation. Hash collisions are excluded
by A4. \qed

\paragraph{Empirical anchor.} The on-chain program v1.4.10 deployed at
\texttt{AmY1ysgQyCC6CmorXkrNkogBSHmomy4kbsPziAYUx47u} on Solana devnet
implements this exact check; the CLI's client-side preflight
(the CLI \texttt{tardus devnet refresh}) rejects double-spend attempts before paying
transaction fees, providing defence in depth.

\subsection{Theorem T4 --- Cut-and-Choose Soundness}
\label{sec:thm-cut-and-choose}

\begin{theorem}[Cut-and-Choose Soundness for Refresh]
\label{thm:cut-and-choose}
Under assumptions A1 (ECDLP) and A2 (ROM), the $\kappa$-fold
cut-and-choose refresh protocol (\S4) binds a malicious
user to producing a denomination-preserving and unlinkable replacement coin
with cheating probability at most
\[
  \Pr[\mathsf{cheat}] \leq \frac{1}{\kappa + 1}.
\]
For the default $\kappa = 32$ (\S4.5), this is
$\approx 3.0\%$ per session; a malicious user wishing to extract
$\Delta$-fold value must succeed $\Delta$ times, with probability
$1 / 33^\Delta$.
\end{theorem}

\paragraph{Proof sketch.} The user commits in Round 2 to $\kappa + 1$
candidate trapdoors $(c_1, \ldots, c_{\kappa+1})$, each encoding a
candidate replacement coin. The mint's Round 3 random challenge selects a
subset $S \subset [\kappa+1]$ of size $\kappa$ to reveal. The remaining
candidate $j = [\kappa+1] \setminus S$ is unblinded and signed; the user
recovers a single fresh coin. For the user to cheat (e.g.\ extract more
denominations than they surrendered), all $\kappa$ revealed candidates
must be consistent with the protocol \emph{and} the one unrevealed
candidate must encode the cheat: probability $1/(\kappa+1)$ assuming the
mint's challenge is uniform. The Round 4 reveal-and-verify step
(\S4.5) checks consistency cryptographically; the
Pedersen binding under A3 prevents equivocation between commit and reveal.

\paragraph{Anonymity preservation.} The unrevealed candidate $j$ is
information-theoretically independent of the user's identity (it is a
freshly sampled scalar). The mint sees only the revealed $\kappa$
candidates, which the user has bound by commitment, plus the blinded
issuance transcript of $j$ which by T2 is independent of the underlying
$(Cp_{\mathsf{new}}, \sigma_{\mathsf{new}})$. Hence the new coin is
unlinkable to the old coin in the union of $\Adv_{\mathsf{net}}$ and
$\Adv_{\mathsf{mint},t-1}$.

\subsection{Theorem T5 --- Reshare Correctness}
\label{sec:thm-reshare}

\begin{theorem}[Proactive Reshare Joint-Key Invariance]
\label{thm:reshare}
Under assumptions A1 (ECDLP), A3 (Pedersen binding), and A5 (honest
threshold), the proactive reshare protocol (\S3.7) transforms
the shares $(s_1, \ldots, s_n)$ into fresh shares $(s_1', \ldots, s_n')$
such that:
\begin{enumerate}[leftmargin=2em,itemsep=0pt]
  \item the joint public key $X = x \cdot G$ is preserved (where
        $x = \sum_{i} \lambda_i s_i = \sum_{i} \lambda_i s_i'$);
  \item the new shares are uniformly distributed conditional on $X$;
  \item old shares from epoch $e$ are useless for any operation in epoch
        $e + 1$ or later.
\end{enumerate}
\end{theorem}

\paragraph{Proof sketch.} The reshare uses a zero-polynomial Pedersen-VSS:
each operator distributes shares of $0$ via a fresh degree-$(t-1)$
polynomial $\tilde{f}_i$ with $\tilde{f}_i(0) = 0$. New shares are
$s_j' = s_j + \sum_i \tilde{f}_i(j)$. The dual-commitment broadcast
(\S3.5) lets honest operators detect a cheating
dealer who chose $\tilde{f}_i(0) \neq 0$; the additional check
$\tilde{A}_{i,0} = \mathcal{O}$ (the identity element)
catches this case explicitly (\S3.7). The
joint-key invariance is algebraic: $\sum_j \lambda_j s_j' =
\sum_j \lambda_j s_j + \sum_i \tilde{f}_i(0) = x + 0 = x$. The
uniformity follows from the freshness of $\tilde{f}_i$. The
\emph{usefulness} of old shares for epoch $e+1$ requires solving
$s_j' = s_j + \delta_j$ for $\delta_j$ given only $X$; this is equivalent
to ECDLP. \qed

\paragraph{Empirical anchor.} The crown-jewel
\texttt{sign\_after\_reshare} test in the \texttt{tardus-mint} crate
(test suite) demonstrates that a signature produced by a
$t$-subset of the new shares verifies under the original $X$, while a
signature attempted with one new and one old share fails.

\subsection{Theorem T6 --- Vault Collateral Invariant}
\label{sec:thm-collateral}

\begin{theorem}[Vault Collateral Invariant]
\label{thm:collateral}
Under assumption A6 (Solana finality) and the on-chain handlers'
canonical-PDA validation (\S5.4), at every
finalised state $\sigma$ of the program the vault's lamport balance
satisfies the bookkeeping invariant
\[
  \mathsf{Vault}[\Denom].\mathsf{lamports} \;\geq\;
  \Denom \cdot \bigl(\,
    \#\mathsf{IssuedCoins}[\Denom] - \#\mathsf{NullifierSet} \cap
        \mathsf{NullifiersOf}[\Denom]
  \,\bigr).
\]
\end{theorem}

\paragraph{Proof.} The handler-by-handler accounting is mechanical:
\texttt{Deposit} requires a preceding system\_program::transfer of at least
$\Denom$ lamports to the vault PDA, observed via
\texttt{vault\_account.lamports()} (\S5.3.2);
\texttt{Withdraw} removes exactly $\Denom$ lamports via
\texttt{invoke\_signed} after a successful nullifier insertion
(\S5.3.4); \texttt{Refresh} does not touch the vault
(\S5.3.3, anonymity preservation). The on-chain code
explicitly checks
\texttt{vault\_account.lamports() < denom} and returns
\texttt{ERR\_VAULT\_INSUFFICIENT} if violated. Hence the invariant is
restored after every successful transaction; A6 ensures durability. \qed

\paragraph{Operational implication.} A user can perform a
\emph{collateral audit} at any time by reading the vault account's
lamport balance and the on-chain nullifier set, then comparing against
the off-chain issued-coin database (publicly aggregated by the mint
committee as part of the transparency log). Any deviation is publicly
observable.

\subsection{Theorem T7 --- Sealed-Box Payload Confidentiality}
\label{sec:theorem-T7}

\begin{theorem}[Sealed-Box IND-CCA, v1.7]
\label{thm:T7}
Let $\mathsf{SB}$ be the TARDUS sealed-box AEAD construction of
\Cref{sec:relay-payload}: $\mathsf{Seal}(m, pk_R) = epk \,\|\,
\mathsf{ChaCha20\text{-}Poly1305}_{K,N}(m)$ where
$(K,N) = \mathsf{HKDF\text{-}SHA\text{-}256}(\mathsf{ECDH}(esk, pk_R^{(\mathrm{X25519})}))$.
For any PPT adversary $\mathcal{A}$ playing the IND-CCA game and
controlling the relay operator (i.e. observing all sealed-box wire
bytes and adaptively requesting decryptions of any non-challenge
ciphertext under the recipient's secret key), there exist PPT
adversaries $\mathcal{B}_1, \mathcal{B}_2$ such that
\[
  \mathsf{Adv}^{\mathrm{IND\text{-}CCA}}_{\mathsf{SB}}(\mathcal{A}) \;\leq\;
  \mathsf{Adv}^{\mathrm{IND\text{-}CCA}}_{\mathsf{ChaCha20\text{-}Poly1305}}(\mathcal{B}_1) \;+\;
  \mathsf{Adv}^{\mathrm{ECDLP}}_{X25519}(\mathcal{B}_2) \;+\;
  \frac{q_H}{2^{256}} \;+\; \mathsf{negl}(\lambda),
\]
where $q_H$ is the number of HKDF (i.e.\ SHA-256-based PRF)
queries.
\end{theorem}

\begin{proof}[Proof sketch]
The reduction composes three layers, each modular against the
others:

\begin{enumerate}[leftmargin=2em,itemsep=0pt]
  \item \textbf{ECDH key indistinguishability.}
    Replace the genuine shared secret
    $s = \mathsf{ECDH}(esk, pk_R^{(\mathrm{X25519})}) = esk \cdot pk_R^{(\mathrm{X25519})}$
    with a uniformly random
    $s' \xleftarrow{\$} \{0,1\}^{256}$ (the ECDLP / DDH game on
    Curve25519). An adversary distinguishing genuine from random
    $s$ in the sealed-box context wins the DDH game for X25519;
    the loss is $\mathsf{Adv}^{\mathrm{ECDLP}}_{X25519}(\mathcal{B}_2)$.
  \item \textbf{HKDF PRF.}
    With $s'$ uniform, the HKDF-derived $(K, N)$ pair is
    statistically indistinguishable from uniform up to a
    $q_H / 2^{256}$ collision probability across the
    $\mathcal{A}$'s HKDF queries.
  \item \textbf{ChaCha20-Poly1305 IND-CCA.}
    The challenge ciphertext is then a fresh ChaCha20-Poly1305
    encryption under uniformly random $(K, N)$. Any IND-CCA
    distinguisher on the sealed-box reduces directly to an IND-CCA
    distinguisher on the underlying AEAD; the loss is
    $\mathsf{Adv}^{\mathrm{IND\text{-}CCA}}_{\mathsf{ChaCha20\text{-}Poly1305}}(\mathcal{B}_1)$.
\end{enumerate}

The compromised-relay adversary model is captured by giving
$\mathcal{A}$ the oracle that returns sealed-box wire bytes for
arbitrary plaintexts (chosen-plaintext oracle) plus the
decryption oracle for any non-challenge ciphertext (chosen-
ciphertext oracle). Both are simulated against the underlying
ChaCha20-Poly1305 oracle in step 3 after the ECDH-to-random hop.

The full mechanisation is in
\texttt{spec/easycrypt/sealed\_box.ec} (Phase 4 audit firm
hand-off).
\end{proof}

\subsection{Theorem T8 --- Confidential Vault Indistinguishability (v1.7 placeholder)}
\label{sec:theorem-T8}

\begin{theorem}[Vault Balance Indistinguishability, v1.5+ target]
\label{thm:T8}
Let $\mathsf{Vault}^{\mathrm{conf}}$ be the v1.5+
Token-2022 ConfidentialMint vault construction (deferred; design
doc at \texttt{deploy/runbooks/token-2022-confidential-mint-design.md}).
For any PPT adversary $\mathcal{A}$ observing the on-chain
encrypted balance commitments + auditor public key + all public
TARDUS metadata, the probability $\mathcal{A}$ distinguishes the
real vault SOL balance $a_0$ from any
$a_1 \in [0, 2^{64})$ is bounded by:
\[
  \mathsf{Adv}^{\mathrm{IND}}_{\mathsf{Vault}^{\mathrm{conf}}}(\mathcal{A}) \;\leq\;
  \mathsf{Adv}^{\mathrm{IND\text{-}CPA}}_{\mathsf{ElGamal}_{\mathrm{Edwards}}}(\mathcal{B}_3) \;+\;
  \mathsf{Adv}^{\mathrm{Soundness}}_{\mathsf{RangeProof}}(\mathcal{B}_4) \;+\;
  \mathsf{negl}(\lambda).
\]
\end{theorem}

\paragraph{Status.} \textbf{Placeholder; not provable on v1
plain-SOL vault.} v1 ships with publicly visible vault SOL
balances (the deposit and withdraw amounts are observable on
the Solana chain). T8 becomes provable once Faz E full ships
(estimated v1.5+, $\approx 19$ engineering weeks; see
architecture decision record in
\texttt{token-2022-confidential-mint-design.md}). The placeholder
is included here to lock the audit firm's expected formal-
verification surface area pre-engagement, so the Phase 4 audit
contract can reference T8 by name.

\paragraph{Reduction sketch (for v1.5+).}
The reduction composes Token-2022's
ElGamal-encrypted-balance ciphertext (the $\mathcal{B}_3$ hop)
with the deposit / withdraw range-proof soundness
($\mathcal{B}_4$ hop). Concretely: the auditor's view sees
$\mathsf{Enc}_{\mathrm{auditor}}(a)$; an IND-CPA distinguisher on
ElGamal under the Edwards subgroup recovers $a$ vs $a'$ with the
$\mathcal{B}_3$ advantage. The range proof's soundness ensures
the amount is in the declared range $[0, 2^{64})$; without
soundness, an adversary could mint negative amounts and break
T6 simultaneously.

\subsection{Theorem T9 --- Fee-Payer Unlinkability (v1.8)}
\label{sec:theorem-T9}

\begin{theorem}[Fee-Payer Unlinkability, Layer 1--4 composed]
\label{thm:T9}
Let $\mathcal{P} = \{P_1, \ldots, P_n\}$ be the set of TARDUS
spend-class transactions (Refresh or Withdraw) submitted to
Solana within an observation window $W$, each accompanied by a
funding chain through the on-chain $\mathsf{SponsorPool}$ PDA per
\Cref{sec:mint-fee-payer-privacy} (Layer 3+4). Let
$\mathcal{F} = \{F_1, \ldots, F_m\}$ be the set of distinct
deposit-class transactions to the $\mathsf{SponsorPool}$ inside
the same window $W$, originating from independent funding
wallets.

For any PPT chain-graph adversary $\mathcal{A}$ with read access
to the entire Solana ledger state inside $W$ — including every
TX hash, every account balance trajectory, every instruction
log, and the public IP source of each TX submission — the
advantage of $\mathcal{A}$ in correctly mapping any specific
spend-class TX $P_i$ to the funding wallet whose deposit
ultimately funded its fee payment is bounded by:
\[
  \mathsf{Adv}^{\mathrm{link}}_{\mathsf{Layer1{+}4}}(\mathcal{A}) \;\leq\;
  \frac{1}{|\mathcal{F}|} \;+\;
  \mathsf{Adv}^{\mathrm{IP\text{-}deanon}}_{\mathrm{network}}(\mathcal{B}_5) \;+\;
  \mathsf{Adv}^{\mathrm{timing}}_{\mathrm{slot}}(\mathcal{B}_6) \;+\;
  \mathsf{negl}(\lambda),
\]
where $\mathcal{B}_5$ wins an IP-deanonymisation game (e.g.\
linking submitter IPs across the Solana RPC layer) and
$\mathcal{B}_6$ exploits per-slot timing fingerprints across the
deposit and payout TXs.
\end{theorem}

\begin{proof}[Proof sketch]
The reduction has three layers:

\begin{enumerate}[leftmargin=2em,itemsep=0pt]
  \item \textbf{Pool indistinguishability ($1/|\mathcal{F}|$ term).}
    The $\mathsf{SponsorPool}$ PDA is a single program-owned account
    whose balance trajectory aggregates all deposits and all
    payouts. From the chain observer's view, every deposit
    $F_j$ contributes its lamports indistinguishably to the
    pool's running balance; every payout $\mathsf{SponsorPayout}$
    instruction (\Cref{sec:onchain-instr-sponsor}) drains a
    fixed amount $\Lambda$ to an ephemeral keypair $E_k$. Inside
    the on-chain log, $E_k$ has no traceable predecessor beyond
    "drawn from the pool"; the funding wallet $F_j$ that
    notionally "covered" the payout is statistically uniform
    over $\mathcal{F}$ as long as deposits arrive faster than
    payouts (a steady-state liquidity assumption).

  \item \textbf{Ephemeral signer + pool caller rotation (Layers 1, 2, 4).}
    Each spend-class TX $P_i$ is signed by a fresh ephemeral
    keypair $E_i$ (Faz 9.1) funded by the pool payout. Each
    pool payout itself is invoked by a caller keypair drawn at
    random from the sponsor pool (Faz 9.5). The on-chain
    signing-chain graph
    $E_i \leftarrow \mathsf{Payout}_{C_j} \leftarrow F_k$
    has fan-in $|\mathcal{F}|$ from every spend's perspective,
    so the joint posterior $\Pr[E_i \text{ funded by } F_k \mid
    \text{public chain state}]$ is $\approx 1/|\mathcal{F}|$.

  \item \textbf{Side channels ($\mathcal{B}_5$, $\mathcal{B}_6$).}
    The reduction does not close (i) network-layer
    deanonymisation (the submitter's IP visible to the Solana
    RPC endpoint and any intermediate router), nor (ii) per-slot
    timing fingerprints (a deposit at slot $S$ and a payout at
    slot $S+\Delta$ for small $\Delta$ creates a correlation
    even if no direct funding link is on-chain). These are
    explicitly enumerated in
    \Cref{sec:mint-fee-payer-privacy} ("Residual leaks not
    addressed by Layers 1--4") and require independent
    countermeasures (Tor / I2P for $\mathcal{B}_5$; random
    submission delays + Jito bundle ordering of
    \Cref{sec:mint-batch} for $\mathcal{B}_6$).
\end{enumerate}

The deposit-rate-exceeds-payout-rate assumption is operationally
necessary: if the pool drains faster than it refills, the
adversary can pinpoint the specific deposit that "must have"
funded a specific payout by lamport-balance accounting. The
spec's drain-rate ceiling
($\approx 14\,\mathrm{SOL/hour}$ at default parameters; see
\Cref{sec:onchain-instr-sponsor}) is calibrated so this
condition is satisfied under typical mainnet deposit cadence.
\end{proof}

\paragraph{Status.} v1.8 spec — theorem statement and informal
reduction shipped. Full mechanisation in EasyCrypt is a Phase 4
audit firm task (the $1/|\mathcal{F}|$ uniform-posterior bound
is straightforward but the steady-state liquidity assumption
requires a real-world model that the audit firm scopes).

\paragraph{Concrete empirical bound (devnet).} The devnet
$\mathsf{SponsorPool}$ PDA holds $\approx 49\,000\,000$ lamports
across $\approx 5$ historical deposits, with payouts of
$10^6$ lamports each. With $|\mathcal{F}| = 5$, an adversary's
direct link probability is $\approx 20\%$ — improving to $< 1\%$
once $|\mathcal{F}| \geq 100$ on mainnet under healthy deposit
cadence. The live demo of 2026-05-22 (TX
\texttt{3UVkfK5tt671FDJ\ldots}) exhibits this property
end-to-end: the spend-class Withdraw signer
\texttt{B3mFCD2qdedYd2nYuTcwZFZS1ghCs7DwKQRtgrvhWszq} is
demonstrably an ephemeral keypair distinct from the deployer
wallet, funded via an on-chain $\mathsf{SponsorPayout}$.

\subsection{Composition: TARDUS-Level Security}
\label{sec:composition}

\begin{corollary}[End-to-End Security]
\label{cor:end-to-end}
The TARDUS protocol satisfies the following composite guarantees:
\begin{enumerate}[leftmargin=2em,itemsep=0pt]
  \item \textbf{Coin authenticity} (T1 + T5): no party can produce a valid
        $(Cp, \sigma)$ without participating in a legitimate issuance or
        refresh session, even across reshare epochs;
  \item \textbf{Sender/receiver/amount privacy} (T2 + T4): on-chain
        observers cannot link a refreshed coin to its predecessor, the
        mint cannot link an issued coin to a refreshed predecessor, and
        the per-coin denomination is the only amount-related leak;
  \item \textbf{Atomic ledger integrity} (T3 + T6): the on-chain vault
        balance and nullifier set are consistent at every finalised
        slot, and double-spends are deterministically rejected.
\end{enumerate}
\end{corollary}

\paragraph{Remaining gaps.} The denomination is leaked at deposit and
withdraw time (the user surrenders exactly $\Denom$ lamports to the vault
PDA, observable on-chain). Phase 1.4.11 will replace the SOL vault with
a Solana Token-2022 Confidential Mint vault, which encrypts the deposit
and withdraw amounts; the per-coin denomination then becomes a private
field of the mint's bookkeeping, not visible on the chain. Phase 1.4.12
will migrate the Borsh-encoded nullifier set
(\S5.3.3) to a Light Protocol compressed Merkle
tree, raising the per-vault capacity from $\approx 250$ to $2^{32}$
nullifiers.

\subsection{Mechanisation Roadmap}
\label{sec:mech-roadmap}

Phase 4 of the project will mechanise theorems T1, T2, and T4 in
EasyCrypt (\cite{easycrypt2013}). Tactical sketch:

\begin{description}[leftmargin=2em,style=nextline]
  \item[\texttt{schnorr.ec}]
    Encodes the single-key Schnorr signature, the
    Pointcheval-Stern forking argument, and the ROM tape. Discharges T1
    in the single-key case.
  \item[\texttt{threshold\_lift.ec}]
    Encodes the simulator for the corrupted minority's view (per
    \cite{stinson-strobl2001}) and reduces threshold T1 to single-key T1.
  \item[\texttt{blind.ec}]
    Encodes the blinding-factor algebra of T2 and proves perfect
    indistinguishability (zero statistical distance) of the issuance
    transcripts.
  \item[\texttt{cut\_choose.ec}]
    Encodes the $\kappa$-fold cut-and-choose game and proves the
    $1/(\kappa+1)$ cheating bound of T4.
\end{description}

T3, T5, and T6 are operational invariants whose mechanisation reduces to
runtime checks already present in the reference implementation; they
require no separate proof effort beyond the test suite of
test suite.
